\chapter{Problem Statement}
\setcounter{equation}{0}
\section{Problem Addressed}
Phishing URL detectors may achieve strong results when training and testing use the same dataset, yet perform less reliably on URLs from other sources. This limits confidence in deploying such models as browser protection. URL features can also vary in availability and predictive value across datasets, making feature selection a key part of generalisation.

\noindent\textbf{Problem statement:} How can an XGBoost-based URL phishing detector be designed and evaluated to balance detection quality, computational efficiency, and generalisation across distinct URL datasets?

\section{Research Questions}
The paper investigates three questions: (1) whether XGBoost achieves cross-dataset generalisation compared with baseline methods; (2) how feature-importance stability relates to cross-dataset performance; and (3) how the computational efficiency of XGBoost compares with neural networks for real-time deployment.\cite{bib1}

\section{Scope of the Seminar}
The scope of this seminar includes the following aspects of the published X-PHIDE method:
\begin{itemize}
\item The motivation for URL-only phishing detection and the challenges caused by dataset shift and changing URL distributions.
\item The sources, composition, and role of the custom, GramBeddings, and PhiUSIIL datasets used in the study.
\item The extraction of lexical, structural, domain-related, and other URL-based features, together with the criteria used to select stable features across datasets.
\item The XGBoost classification pipeline, Bayesian hyperparameter optimisation, decision thresholds, and comparison with baseline and neural-network models.
\item The evaluation measures used to assess classification quality, cross-dataset generalisation, training and inference time, and memory requirements.
\item The proposed browser plug-in workflow, including allowlist and blocklist checks and the presentation of URL risk categories to users.
\item The reported strengths, limitations, practical deployment considerations, and possible directions for future research.
\end{itemize}

The report focuses on explaining and analysing the methods and results presented by Misiek and Hyla. It does not include independent data collection, model training, statistical re-analysis, or implementation of a complete production-ready browser extension. All numerical results and performance claims are attributed to the original authors; no independent reproduction is claimed.\cite{bib1}
